\documentclass[10pt,a4paper]{amsart}
\usepackage[margin=19mm]{geometry}
\usepackage{amsmath,amssymb,mathtools}
\usepackage[scr=rsfs,cal=euler]{mathalfa}
\usepackage{xfrac}
\usepackage[unicode,hidelinks,bookmarks=false]{hyperref}
\newcommand{\ie}{i.e.\ }
\newcommand{\cf}{cf.\ }
\newcommand{\resp}{respectively\ }
\newcommand{\Subsection}{\subsection}
\providecommand{\nobreakdash}{\nobreak\hbox{-}}
\setlength{\emergencystretch}{2em}
\multlinegap0pt
\usepackage{xcolor}
\usepackage{amssymb}
\usepackage{algorithm}\usepackage{algpseudocode}
\usepackage{tikz}
\usepackage{caption}
\usepackage{natbib}
\newtheorem{lemma}{Lemma}
\newcommand{\FRil}[1]{\todo[color=black, textcolor=white, inline, author=FR]{#1}}
\newcommand{\X}{{\mathcal{X}}} % Solution space
\newcommand{\fsel}{S} % Set of selected instance features
\newcommand{\instanceindex}{i}
\newcommand{\instanceindextwo}{j}
\newcommand{\instfeatures}{[p]} % Instance Features
\newcommand{\neighbour}{\ensuremath{\mathcal{N}}}
\newcommand{\sEpsBorder}[1][(\instanceindex,\selectionset)]{\ensuremath{\neighbour^=_\epsilon#1}}
\newcommand{\sEpsNeighbour}[1][(\instanceindex,\selectionset)]{\ensuremath{\neighbour^\le_\epsilon#1}}
\newcommand{\sEpsStrict}[1][(\instanceindex,\selectionset)]{\ensuremath{\neighbour^<_\epsilon#1}}
\newcommand{\selectionset}{\fsel}
\newcommand{\skPess}[1][(\instanceindex,\selectionset)]{\neighbour_k^\textnormal{pess}#1}
\newcommand{\soldist}[1][(\instanceindex,\instanceindextwo)]{d_\X#1} % Solution distance
\newcommand{\solution}{x}
\title{Feature Selection for Data-Driven Explainable Optimization}
\author{Kevin-Martin Aigner, Marc Goerigk, Michael Hartisch, Frauke Liers, Arthur Miehlich, Florian R\"osel}
\date{Source: arXiv:2504.12184v2 (CC BY 4.0)}
\begin{document}
\maketitle

\noindent\emph{Source and licence.} Feature Selection for Data-Driven Explainable Optimization. Version arXiv:2504.12184v2 (CC BY 4.0), distributed by arXiv under the Creative Commons Attribution 4.0 International licence, http://creativecommons.org/licenses/by/4.0/, as recorded in the arXiv licence metadata for this exact version and shown on the abstract page https://arxiv.org/abs/2504.12184v2. Full licence notice: \texttt{licenses/arxiv-2504.12184-CC-BY-4.0.txt}.\par\smallskip

\noindent\textit{Editorial scope.} Counted unit: the article's lemma lemma:bound\_on\_alpha (source0.tex lines 1309--1314). The article's complete original proof of it is delivered with it (source0.tex lines 1316--1338, one \texttt{proof} environment of 23 lines). No mathematics is written, repaired or re-derived: the statement and the proof are the article's own lines, in the article's own notation, and no retained line is substituted or re-flowed.  Retained uncounted as the source-local context the proof needs: lines 1200--1210; lines 1243--1275; lines 1277--1300.  Nothing was omitted from any retained span, so the package carries no omission record.  No retained line contains a citation, so the wrapper carries no bibliography and no bibliography entry.  No retained line contains a figure environment, an image-inclusion command or drawing code, so no image is delivered and the package declares no asset dependency.  Editorial formatting only, in addition: an AMS \texttt{amsart} preamble that restates the article's own declarations and macros used by the retained lines, namely the environments \texttt{lemma} on separate counters, together with the standard packages those lines need.  The retained environments print in this document as Lemma 1 in order of appearance, whereas the article prints the same items with the numbering of its own full document.  The complete native source of the article as distributed in the arXiv e-print for this exact version is bundled unchanged as \texttt{sources/176652/source0.tex}.
	\begin{align}
		\label{problem:optimistic_fspeo:obj}\min \quad &\sum_{i=1}^N \sum_{j=1,j\neq i}^N\ y_{ij}\cdot \soldist\\
		\label{problem:optimistic_fspeo:k_neighbors}\text{s.t.} \quad &\sum_{j=1, j\neq i}^N y_{ij} \ge k\quad \forall i\in [N],\\
		\label{problem:optimistic_fspeo:L_features}&1 \le \sum_{f\in\instfeatures} b_f\le L,\\
		\label{problem:optimistic_fspeo:instance_diff}&d_{ij} = \sum_{f\in\instfeatures} b_f\ d_f(i,j) \quad \forall i\neq j\in [N],\\
		\label{problem:optimistic_fspeo:neighboring1}&d_{ij} \le\epsilon_i+(1-y_{ij})M\quad \forall i\neq j\in [N],\\
		\label{problem:optimistic_fspeo:neighboring2}&\epsilon_i\le d_{ij} + y_{ij}M\quad \forall i\neq j\in [N],\\
		\label{problem:optimistic_fspeo:vardef}&b\in\{0,1\}^{p},\; y\in\{0,1\}^{N\times N-1},\; d \in \mathbb{R}^{N\times N-1}_{\geq 0},\;
		%&y_{ii} = 0 \quad \forall i \in [N]\\
	    \epsilon\in\mathbb{R}_{\ge 0}^N.
	\end{align}

We first study the evaluation problem for the case that the instance distances $d_{ij}$ are fixed to some positive values.
The evaluation problem then reads:
\begin{subequations}\label{problem:eval}
	\begin{align}
		\max \quad & \sum_{i=1}^N \sum_{j=1, j\neq i}^Ny_{ij} \cdot \soldist \\
		\label{problem:eval_bound}\text{s.t.}\quad & \sum_{j=1, j\neq i}^N d_{ij}y_{ij} \leq \min \Big\{ \sum_{j=1, j\neq i}^N d_{ij}y^\prime_{j} \colon \sum_{j=1, j\neq i}^N y^\prime_j \geq k, y^\prime \in [0,1]^{N-1} \Big\} \quad \forall i\in [N],\\
		\label{problem:eval_k}& \sum_{j=1,j\neq i}^N y_{ij} \leq k \quad \forall i\in [N],\\
		%& d_{ij} = \sum_{l\in[F_I]} b_l\ d_l(I^i,I^j) \quad \forall i,j\in [N]\\
		\label{problem:eval_integrality}& y_{ij}\in \{0, 1\}^{N \times N-1}.
	\end{align}
\end{subequations}
% \FRil{$d_\X$...}
The binary variables $y_{ij}$ are equal to $1$ if instance $I^j$ is in $\skPess$.  
We want to note that Problem~\eqref{problem:eval} decomposes into $N$ smaller optimization problems, one for each $i \in [N]$. Each of these is a cardinality constrained knapsack problem.
\begin{lemma}
	The integrality constraints \eqref{problem:eval_integrality} are redundant as the continuous relaxation of \eqref{problem:eval} has always an integer solution.
\end{lemma}
%   \textcolor{red}{FL:den folgenden Satz fand ich schwer verstaendlich, habe
%     umformuliert - schaust du mal bitte, Florian?}
%     \FRil{Lass mal drüber reden.}
\begin{proof}
We show that vectors with fractional $y_{ij}$ entries do not
  correspond to vertices of the linear programming relaxation.
  %we can only use features with fractional
   %     variable values from those with equal weight on the boundary of the $k$-th smallest:
	Let $\epsilon > 0$ be such that $ |\sEpsStrict| \leq k \leq |\sEpsNeighbour|$. Then the right hand side of Constraint~\eqref{problem:eval_bound} evaluates to
	\begin{equation*}
		\sum_{j \in \sEpsStrict} d_{ij} + (k - \lvert \sEpsStrict \rvert) \epsilon.
	\end{equation*}
	Hence, the variables $y_{ij}$ have value $1$ for all $j \in \sEpsStrict$, and have value $0$ for all $j\notin \sEpsNeighbour$.
	If further a variable $y_{ij}$ with $j \in \sEpsBorder$ has a fractional value, another index $j^\prime \in \sEpsBorder$ exists with $y_{ij^\prime}$ fractional due to Constraint~\eqref{problem:eval_k}. As $d_{ij} = d_{ij^\prime}$, increasing $y_{ij}$ by $\min\{(1 - y_{ij}), y_{ij^\prime}\}$ and decreasing $y_{ij^\prime}$ by the same number preserves feasibility. Hence, no fractional $y$ can be a vertex of the feasible set of the continuous relaxation of Problem~\eqref{problem:eval} and the integrality constraints can be removed without consequences.
%	\Halmos
\end{proof}

We dualize the inner minimization problem to replace Constraint~\eqref{problem:eval_bound} by
\begin{subequations}\label{problem:eval_dualize_inner}
	\begin{align}
		%\max \quad & \sum_i \sum_j d_\X(\phi_\X(x^i),\phi_\X(x^j)) y_{ij}\\
		& \sum_{j=1, j\neq i}^N d_{ij}y_{ij} \leq k \pi_i - \sum_{j=1, j\neq i}^N \rho_{ij} \quad \forall i\in [N],\\
		&\pi_i - \rho_{ij} \leq d_{ij} \quad \forall i\neq j\in [N],\\
		%& \sum_j y_{ij} = k \quad \text{for all }i\\
		%& d_{ij} = \sum_{l\in[F_I]} b_l\ d_l(I^i,I^j) \quad \forall i,j\in [N]\\
		%& y_{ij}\in [0, 1] \quad \forall i,j\in [N] \\
		& \pi_i \geq 0, \rho_{ij} \geq 0 \quad \forall i\neq j\in [N].
	\end{align}
\end{subequations}
and dualize Problem~\eqref{problem:eval}, considering Constraint~\eqref{problem:eval_k} as inequality, and express the $d_{ij}$ as a combination of at most $L$ instance features to get
\begin{subequations}\label{problem:eval_dualize_outer}
	\begin{align}
		\min \quad & \sum_{i=1}^N \sum_{j=1,j\neq i}^N d_{ij} \beta_{ij} + \sum_{i=1}^N k \gamma_i + \sum_{i=1}^N \sum_{j=1, j\neq i}^N \delta_{ij}\\
		\nonumber\text{s.t.} \quad & \eqref{problem:optimistic_fspeo:L_features}, \eqref{problem:optimistic_fspeo:instance_diff}\\
		\label{problem:eval_dualize_outer:b_cons}& d_{ij} \alpha_i + \gamma_i + \delta_{ij} \geq \soldist\quad \forall i\neq j\in [N],\\
		\label{problem:eval_dualize_outer:alpha_bound}& \sum_{j=1,j\neq i}^N \beta_{ij} \geq k \alpha_i \quad \forall i\in [N],\\
		\label{problem:eval_dualize_outer:beta_bound}& \alpha_i \geq \beta_{ij} \quad \forall i\neq j\in [N],\\
		&\alpha, \beta, \gamma, \delta, d \geq 0, \;b \in \{0, 1\}^{p}.
	\end{align}
\end{subequations}
% \FRil{$d_\X$...}

\begin{lemma}\label{lemma:bound_on_alpha}
	There is an optimal solution $(\alpha^*, \beta^*, \gamma^*, \delta^*, d^*, b^*)$ to Problem~\eqref{problem:eval_dualize_outer} that fulfills
	\begin{equation}
		\label{alphabound}\alpha_i^* \leq \max_{j \in [N]\setminus\{i\}} \frac{\soldist}{\min_{f\in \instfeatures: d_f(i,j) \neq 0} d_f(i,j)}
	\end{equation}
\end{lemma}

\begin{proof}
	Assume that Inequality~\eqref{alphabound} does not hold for  $i \in [N]$. Then we can construct another optimal solution by
        setting $\tilde{\alpha}_i^*$ to the right-hand side of the inequality, and $\tilde{\beta}^*_{ij}$ to $\frac{\beta^*_{ij}\tilde{\alpha}_i^*}{\alpha_i^*}$ for all $j\in [N]$. For this, we consider two sets $M^0 := \{j \in [N]\setminus \{i\} \colon d^*_{ij} = 0\}$ and $M^+:= \{j \in [N]\setminus \{i\} \colon d^*_{ij} \neq 0\}$. For $j \in M^0$, it holds that
	\begin{equation*}
		\gamma_i^* + \delta^*_{ij} \geq d_\X(\solution^i,\solution^j),
	\end{equation*}
	hence this constraint does not restrict our choice on $\alpha_i$. For $j \in M^+$ it holds that
	\begin{align*}
		d_{ij}^* \tilde{\alpha}_i + \gamma_i^* + \delta_{ij}^*& \geq \min_{f\in \instfeatures: d_f(i,j) \neq 0}d_f(i,j) \frac{\soldist}{\min_{l\in \instfeatures: d_f(i,j) \neq 0} d_f(i,j)} + \gamma_i^* + \delta_{ij}^*\\ 
		&\geq \min_{f\in \instfeatures: d_l(i,j) \neq 0}d_f(i,j) \frac{\soldist}{\min_{f\in \instfeatures: d_f(i,j) \neq 0} d_f(i,j)}\\ 
		&= \soldist. %d_\X(\solution^i,\solution^j).
	\end{align*}
Since
	\begin{equation*}
		\sum_{j =1,j\neq i}^N \tilde{\beta}_{ij}^* = \sum_{j =1,j\neq i}^N \frac{\beta_{ij}^* \tilde{\alpha}^*_i}{\alpha_i^*} = \frac{\tilde{\alpha}^*_i}{\alpha_i^*} \sum_{j =1,j\neq i}^N \beta_{ij}^* \geq \frac{\tilde{\alpha}^*_i}{\alpha_i^*} k \alpha_i^* = k \tilde{\alpha}_i^*
	\end{equation*}
the new solution fulfills Constraint~\eqref{problem:eval_dualize_outer:alpha_bound}.
We calculate that with
	\begin{equation*}
		\tilde{\alpha}_i^* = \frac{\tilde{\alpha}_i^* \alpha_i^*}{\alpha_i^*} \leq \frac{\tilde{\alpha}_i^* \beta_{ij}^*}{\alpha_i^*} = \tilde{\beta}_{ij}^*,
	\end{equation*}
Constraint~\eqref{problem:eval_dualize_outer:beta_bound} is also satisfied. Hence the solution we constructed is feasible and fulfills Inequality~\eqref{alphabound}. As $\tilde{\beta}_{ij}^* \leq \beta_{ij}^*$ for all $j \in [N]\setminus \{i\}$, its objective value is not worse than the value of the original solution. This proves the claim. %\Halmos
\end{proof}

\end{document}
